Pure functional languages that compile to JavaScript must choose how to represent their immutable
sequences. A persistent structure makes every update cheap and every version safe to keep, but it
costs runtime code, indirection on every read, and a representation that JavaScript cannot consume
directly. A plain JavaScript array written in place is what a hand-written program uses, but it is
correct only when no one can observe the old version. This thesis designs an \emph{inferred
ownership checker} for beni, a strict, pure, Elm-like language for the browser, under which an
in-place edit of a list while anything else still holds the old version is a \emph{compile-time
error}, with no ownership annotations in user code. Every accepted in-place edit then writes its
array, in development and release builds alike, with no reference counts and no silent copy. List
literals, spreads and concatenation keep their meaning: they build a new array, as a JavaScript
array literal does, and in-place edits are spelled with the library's named writers.

The property checked is uniqueness \emph{at the update site}, decided by liveness: every other
holder of the array must be dead once the update's arguments have been evaluated. A dereference that
completes before the write is therefore never sharing, and a dead alias is harmless. Strict
left-to-right evaluation fixes when each read happens, saturated calls rule out silent captures by
partial application, and the type checker already infers per-function facts and publishes them in
module interfaces. The analysis is a directional, flow-sensitive, field-sensitive abstraction of
which arrays each value may hold---including the environments of closures, so that a closure
returned, stored or captured carries what it holds and whether it may be called only once---a
per-path liveness pass, and four side conditions at each update. Function summaries, polymorphic
over the ownership behaviour of function-typed parameters, are solved for recursive groups by an
optimistic fixpoint over a finite lattice and published in interfaces. An \emph{exclusive-contents}
flag, meaning that no two list positions inside a container share an array, extends the rule to
edits inside the elements of a container.

We present the model, the inference algorithm and its complexity, a sketch of a soundness argument
(the value semantics and the in-place semantics produce the same observable trace) under twelve
explicitly stated assumptions, the trusted surface those assumptions cover, an entry protocol for
every function a UI runtime calls, a treatment of the hard cases, an error-message design that
separates real sharing from the checker's own limits, the changes to the language and its library
that the rule needs or invites, and worked programs, including twelve that an earlier version of
the rules wrongly accepted and that the rules now reject. The design rejects some correct programs in
exchange for a guarantee about cost; we state that cost plainly. Two applications we examined
contain no in-place edit under the rule and gain nothing from it as written. We propose a
report-mode experiment, with a differential run that traps every read of a stale array, before the
rule is adopted. No results of that experiment exist yet, and the proof is a sketch that has not
been mechanised.
